\section{Discussion}
\subsection{When traditional ML is competitive}
The Java results demonstrate that an encoder is not automatically preferable for these partitions. Word/character TF-IDF with logistic regression matches E1 screening micro-F1 to within a few hundredths of a percentage point and has higher macro-F1. Classical final results also remain close to the neural leaders, with the chain producing the highest observed Java macro-F1 among available final methods. This supports retaining lexical baselines in model selection. It does not establish deployment efficiency: the repository lacks a controlled comparison of inference latency, memory, and hardware-normalized energy or cost.

The larger reference-model advantage on Pharo suggests that the balance between lexical cues and contextual representation differs across the corpora. Different taxonomies, text distributions, and support profiles vary together with programming language. It is therefore not possible to isolate programming language as the cause. The present evidence supports a corpus-specific comparison, not a general ranking of how well Transformers understand the three languages.

\subsection{Backbones, objectives, and interactions}
The experiment matrix contains useful contrasts, including E1/E11 for backbone changes and E11/E15 for CodeBERT R-Drop. However, E15 versus E1 changes both backbone and consistency regularization, and E16 combines learning-rate changes with R-Drop. E9 also has a different stopping policy. A headline comparison between each language's strongest configuration and E1 cannot allocate gains uniquely among backbone, objective, stopping, and optimization effects. No balanced factorial study or controlled causal analysis is present.

The observation that CodeBERT leads screening on Python but its BCE variant trails RoBERTa on Pharo is compatible with several hypotheses: differences in lexical cues, pretraining relevance, head fitting, and label-support interactions. None is verified by the available artifacts. The report therefore separates observed score differences from proposed explanations. Attention weights, graph design, or model complexity alone do not establish interpretable reasoning about comments.

\subsection{Aggregate and rare-label performance}
Sparse positive assignments create different objectives for different use cases. A documentation-search tool prioritizing common usage information may reasonably emphasize micro-F1. A maintenance tool requiring reliable detection of DevelopmentNotes, deprecation, or Classreferences needs label-specific precision and recall. Macro-F1 gives each category equal weight but can still mask a particular regression by averaging improvements elsewhere. The explicit Python DevelopmentNotes regression under E15 is a reason to inspect the relevant label, even when both aggregate F1 measures increase.

Changing a threshold can alter this tradeoff without changing the encoder. The fixed 0.5 threshold is an operational choice, not a demonstrated optimum for every label. E4 implements validation-based tuning but was not evaluated in the final cohort. Its unchanged screening scores should not be used to dismiss threshold calibration. A future application would need validation criteria reflecting the actual costs of false positives and false negatives.

\subsection{Practical implications and generalization}
A defensible workflow would keep a simple lexical baseline, compare neural candidates using training-only selection, inspect label-level outcomes, and evaluate a frozen configuration on genuinely unseen groups. The current repository already supplies shared data validation and transparent result contracts that would support that extension. The new dataset adapters extend its storage design; a deployed inference service and independent project-level evaluation remain future work.

The appropriate practical claim is that these methods can learn information-type associations on this dataset under the recorded protocol. Reliability on new projects, different comment conventions, other languages, and developer workflows is unresolved. Potential downstream gains in maintenance productivity or documentation quality require their own evaluation.
